\subsection{The relative separable algebraic closure}

PROPOSITION 12. --- Let E be an extension of K and \(E_{s}\) the set of elements of E which are algebraic and separable over K. Then \(E_{s}\) is the largest subextension of E which is algebraic and separable over K.

By Prop. 6, a) (V, p.~39) every subextension of E which is algebraic and separable over K is contained in \(E_{s}\) . By Prop. 6, b) (loc. cit.) the extension \(\mathrm{K}(E_{s})\) of K is algebraic and separable, whence \(\mathrm{K}(E_{s}) \subset E_{s}\) and so \(\mathrm{K}(E_{s}) = E_{s}\) .

With the notation of the preceding proposition, \(E_{s}\) is called the relative separable (algebraic) closure of K in E. When K is perfect, \(E_{s}\) is the relative algebraic closure of K in E (V, p.~36, Prop. 2).

PROPOSITION 13. --- Let E be an algebraic extension of K and let \(E_{s}\) be the relative separable algebraic closure of K in E.

\begin{enumerate}
\def\labelenumi{\alph{enumi})}
\item
  E is a \(p\) -radical extension of \(\mathrm{E}_s\) .
\item
  If \(\mathbf{F}\) is a subextension of \(\mathbf{E}\) such that \(\mathbf{E}\) is \(p\) -radical over \(\mathbf{F}\) , then \(\mathbf{F} \supseteq \mathbf{E}_s\) .
\item
  \(\mathbf{E}_s\) is the unique subextension of \(\mathbf{E}\) which is separable over \(\mathbf{K}\) and over which \(\mathbf{E}\) is \(p\) -radical.
\end{enumerate}

It suffices to prove \(a\) ) in the case where \(\mathbf{K}\) is of characteristic \(p \neq 0\) . Let \(x\) be an element of \(\mathbf{E}\) and \(f\) its minimal polynomial over \(\mathbf{K}\) . There exists an integer \(m \geqslant 0\) such that \(f\) belongs to \(\mathbf{K}[\mathbf{X}^{p^m}]\) but not to \(\mathbf{K}[\mathbf{X}^{p^{m+1}}]\) ; in other words, we have \(f(\mathbf{X}) = g(\mathbf{X}^{p^m})\) with \(g \in \mathbf{K}[\mathbf{X}], g \notin \mathbf{K}[\mathbf{X}^p]\) . Since \(f\) is irreducible, the same is true of \(g\) , hence \(g\) is the minimal polynomial of \(x^{p^m}\) over \(\mathbf{K}\) . By \(V\) , p.~38, Prop. 4 and p.~39, Prop. 5 we thus have \(x^{p^m} \in \mathrm{E}_s\) , so \(\mathbf{E}\) is \(p\) -radical over \(\mathrm{E}_s\) .

Assume now the hypothesis \(b)\) and let \(x \in \mathrm{E}_s\) . Since \(x\) is separable over \(\mathbf{K}\) , it is so also over \(\mathrm{F}(\mathrm{V}, \mathrm{p.~39}, \mathrm{Cor.~2})\) , but since \(\mathrm{E}\) is \(p\) -radical over \(\mathrm{F}\) , \(x\) is also \(p\) -radical over \(\mathrm{F}\) , hence \(x \in \mathrm{F}(\mathrm{V}, \mathrm{p.~39}, \mathrm{Cor.~3})\) .

Finally c) follows from a) and b) and Prop. 12.

COROLLARY 1. --- Let E and K' be two extensions of K contained in the same extension of K. Suppose that E is algebraic over K and denote by \(E_{s}\) the relative separable algebraic closure of K in E. Then \(\mathrm{K}'(\mathrm{E}_{s})\) is the relative separable algebraic closure of \(K'\) in \(\mathrm{K}'(\mathrm{E})\) .

For \(\mathbf{K}^{\prime}(\mathbf{E}_{s})\) is a separable algebraic extension of \(K^{\prime}\) by Prop. 10 (V, p.~42); since E is p-radical over \(E_{s}\) , the extension \(\mathbf{K}^{\prime}(\mathbf{E})\) of \(\mathbf{K}^{\prime}(\mathbf{E}_{s})\) is p-radical (V, p.~25, Cor.). Now it suffices to apply Prop. 13.

COROLLARY 2. --- If E has finite degree over K, then \(E_{s} = \bigcap_{n \geqslant 0} K(E^{p^{n}})\) .

For each integer \(n \geqslant 0\) denote by \(F_n\) the subextension \(K(E^{p^n})\) of \(E\) . The sequence \((F_n)_{n \geqslant 0}\) of vector subspaces of \(E\) is descending and \(E\) is of finite dimension over \(K\) . Hence there exists an integer \(m \geqslant 0\) such that \(F_m = F_n\) for all \(n \geqslant m\) . We thus have \(K(F_m^p) = F_{m+1} = F_m\) , hence \(F_m\) is a separable extension of \(K(V, p. 42, \text{Prop. } 11)\) ; it is clear that \(E\) is \(p\) -radical over \(F_m\) and so Prop. 13 implies \(E_s = F_m = \bigcap_{n \geqslant 0} F_n\) .

Remark. --- Let E be an algebraic extension of K and \(E_{r}\) the relative p-radical closure of E in K (V, p.~25). Then \(E_{r}\) is the largest subextension of E which is p-radical over K (V, p.~25, Prop. 2). However, E is in general not separable over \(E_{r}\) (V, p.~152, Ex. 2); for the case of quasi-Galois extensions see V, p.~76.

